%=====================================================================
\chapter{Generative AI and Large Language Models}\label{ch:genai}
%=====================================================================
\locator{4}{Part IV --- Deep Learning and Modern AI}

\begin{outcomes}
By the end of this chapter you will be able to:
\begin{tight}
  \item \textbf{Explain} how an LLM produces text, and why that mechanism causes hallucination.
  \item \textbf{Distinguish} pre-training, fine-tuning, prompting and retrieval, and choose between them.
  \item \textbf{Write} effective prompts using role, context, task, format and examples.
  \item \textbf{Describe} the RAG architecture and say which problem each component solves.
  \item \textbf{Evaluate} a generative system despite the absence of a single correct answer.
\end{tight}
\end{outcomes}

\begin{prereq}
\chref{architectures} (attention and transformers) and \chref{nlp} (embeddings ---
RAG is built on them).
\end{prereq}

\begin{keyterms}
generative model $\bullet$ large language model $\bullet$ token $\bullet$
next-token prediction $\bullet$ temperature $\bullet$ context window $\bullet$
pre-training $\bullet$ fine-tuning $\bullet$ prompt engineering $\bullet$
zero/few-shot $\bullet$ chain of thought $\bullet$ hallucination $\bullet$
retrieval-augmented generation $\bullet$ vector database $\bullet$ grounding
\end{keyterms}

\section{How an LLM Actually Works}

\begin{definitionbox}[Large Language Model]
A very large transformer (\chref{architectures}) trained on an enormous corpus with
one objective: given a sequence of tokens, predict the next one. Text is generated
by doing this repeatedly, each new token appended and fed back in.
\end{definitionbox}

\dsfig{%
\begin{tikzpicture}[font=\scriptsize,
  tok/.style={rounded corners=2pt, draw=primary!60, fill=primary!10, text=primary!55!black,
              minimum width=1.15cm, minimum height=0.5cm, font=\tiny},
  new/.style={rounded corners=2pt, draw=greenhl, fill=greenhl!16, text=greenhl!45!black,
              minimum width=1.15cm, minimum height=0.5cm, font=\tiny\bfseries}]

\foreach \i/\w in {0/The, 1/firewall, 2/blocked, 3/the}{
  \node[tok] (t\i) at (\i*1.35,1.4) {\w};}
\node[new] (t4) at (4*1.35,1.4) {?};

\node[rounded corners=3pt, draw=purplehl, fill=purplehl!8, text=purplehl,
      minimum width=3.0cm, minimum height=0.8cm, font=\tiny\bfseries] (m) at (2.7,0.15)
     {TRANSFORMER};
\foreach \i in {0,...,3}{
  \draw[-{Stealth[length=1.4mm]}, gray!55] (t\i) -- (m);}

% probability distribution over next token
\begin{scope}[shift={(7.2,-0.55)}]
  \node[font=\tiny\bfseries, text=primary, anchor=west] at (0,2.15)
       {probability over the whole vocabulary};
  \foreach \i/\w/\p/\bw in {0/connection/0.34/2.72, 1/request/0.21/1.68,
                            2/traffic/0.18/1.44, 3/attack/0.11/0.88,
                            4/banana/0.0001/0.05}{
    \pgfmathsetmacro{\yy}{1.75-\i*0.42}
    \node[font=\tiny, anchor=east, text=gray!50!black] at (-0.15,\yy) {\w};
    \draw[rounded corners=1pt, draw=secondary, fill=secondary!35]
         (0,\yy-0.13) rectangle (\bw,\yy+0.13);
    \node[font=\tiny, anchor=west, text=gray!50!black] at (\bw+0.12,\yy) {\p};}
\end{scope}
\draw[-{Stealth[length=1.8mm]}, purplehl, line width=1pt] (m) -- (6.3,0.15);

\draw[-{Stealth[length=1.8mm]}, greenhl, line width=1pt]
     (10.6,1.2) .. controls +(0.9,0.9) and +(0.9,-0.9) .. (t4.east);
\node[font=\tiny, text=greenhl!45!black, anchor=west] at (11.3,1.4) {sample one};

\node[anchor=north west, align=left, text width=13.3cm, font=\scriptsize, text=accent!75!black]
     at (-1.0,-2.15)
     {\textbf{This is the whole mechanism, and it explains the failure mode.} The model
      samples a \emph{plausible} next token, not a \emph{true} one. It has no separate
      store of facts to check against and no representation of its own uncertainty about
      the world. A fabricated citation is generated by exactly the same process, with
      exactly the same confidence, as a correct one --- which is why hallucination cannot
      be fixed by asking the model to be more careful.};
\end{tikzpicture}}%
{Next-token generation. The model outputs a probability distribution over its vocabulary
and samples from it; the sampled token is appended and the process repeats. Fluency and
truth are produced by the same mechanism, which is the central risk.}{llm-generation}

\begin{conceptbox}[Temperature]
\term{Temperature} controls how sharply the model samples. Near 0 it always takes
the most likely token --- repetitive but predictable, and the right setting for
extraction, classification and code. Around 0.7--1.0 it samples more freely ---
varied and creative, and the right setting for drafting and brainstorming. Setting
temperature is a task decision, not a style preference.
\end{conceptbox}

\section{Four Ways to Adapt a Model}

\dsfig{%
\begin{tikzpicture}[font=\scriptsize,
  hd/.style={rounded corners=3pt, fill=#1, text=white, font=\scriptsize\bfseries,
             text width=3.35cm, align=center, minimum height=0.55cm},
  b/.style={rounded corners=3pt, draw=#1, fill=#1!8, text width=3.35cm, align=center,
            font=\scriptsize, inner sep=5pt, minimum height=3.3cm}]

\node[hd=primary] at (0,0)    {PROMPTING};
\node[b=primary]  at (0,-2.05){%
  Instructions only; the model is unchanged.\\[4pt]
  \textbf{Cost:} nearly none\\
  \textbf{Data:} none\\
  \textbf{Time:} minutes\\[4pt]
  {\color{greenhl!45!black}Start here. Always.}};

\node[hd=secondary] at (3.9,0)    {RAG};
\node[b=secondary]  at (3.9,-2.05){%
  Retrieve your documents, put them in the prompt.\\[4pt]
  \textbf{Cost:} low\\
  \textbf{Data:} a document store\\
  \textbf{Time:} days\\[4pt]
  {\color{greenhl!45!black}The fix for ``it doesn't know our stuff''.}};

\node[hd=goldhl] at (7.8,0)    {FINE-TUNING};
\node[b=goldhl]  at (7.8,-2.05){%
  Continue training on your labelled examples.\\[4pt]
  \textbf{Cost:} moderate\\
  \textbf{Data:} 1{,}000+ examples\\
  \textbf{Time:} weeks\\[4pt]
  {\color{accent}For \emph{behaviour} and format, not for facts.}};

\node[hd=accent] at (11.7,0)    {PRE-TRAINING};
\node[b=accent]  at (11.7,-2.05){%
  Train a model from nothing.\\[4pt]
  \textbf{Cost:} millions\\
  \textbf{Data:} trillions of tokens\\
  \textbf{Time:} months\\[4pt]
  {\color{accent}Not something you will do.}};

\draw[-{Stealth[length=2.5mm]}, gray!40, line width=5pt, opacity=0.5]
     (-1.9,-4.15) -- (13.6,-4.15);
\node[font=\scriptsize\bfseries, text=gray!55!black] at (5.9,-4.55)
     {increasing cost, data requirement and commitment $\longrightarrow$};

\node[anchor=north west, align=left, text width=13.4cm, font=\scriptsize, text=accent!75!black]
     at (-1.9,-5.35)
     {\textbf{The distinction students most often get wrong:} fine-tuning teaches a model
      \emph{how to behave} --- tone, format, task pattern. It is a poor and expensive way
      to teach it \emph{what is true}, because facts change and retraining does not. For
      knowledge, use RAG.};
\end{tikzpicture}}%
{Four levels of adaptation. Work left to right and stop as soon as the problem is solved;
most production systems never move past the second box.}{llm-adaptation}

\section{Prompting}

\begin{definitionbox}[The Five Components of a Good Prompt]
\term{Role} (who the model should act as), \term{context} (the background it
needs), \term{task} (the specific instruction), \term{format} (the required output
structure), and \term{examples} (one or more demonstrations). Missing components are
the usual cause of a disappointing answer.
\end{definitionbox}

\begin{worked}[Improving a Prompt]
\textbf{Attempt 1 --- vague.}\\
\texttt{Analyse this student data.}\\
\emph{Result:} a generic essay about what data analysis is. No role, no task, no
format.

\smallskip
\textbf{Attempt 2 --- task added.}\\
\texttt{List which of these students are at risk of failing.}\\
\emph{Result:} a plausible-looking list with no stated reasoning and no way to
check it. Better, but unusable and unauditable.

\smallskip
\textbf{Attempt 3 --- all five components.}
\begin{lstlisting}[basicstyle=\ttfamily\tiny, frame=none, backgroundcolor=\color{codebg}]
ROLE:    You are an academic support adviser.
CONTEXT: Below are week-4 engagement records. Students fail mainly through
         disengagement, not ability. We can contact at most 20 students.
TASK:    Identify the 20 students most likely to fail. For each, give the
         single strongest reason, drawn only from the data provided.
FORMAT:  A markdown table: student_id | risk (high/medium) | reason | evidence
RULES:   Use only the supplied data. If evidence is insufficient for a
         student, write "insufficient data" rather than guessing.
EXAMPLE: S101 | high | stopped submitting in week 2 | 0 submissions wk2-4
\end{lstlisting}

\smallskip
\textbf{Why attempt 3 works.} The role sets vocabulary and stance; the context
supplies the constraint (20 students) that makes the task well-posed; the format
makes the output machine-readable and comparable; the rule against guessing gives
the model an explicit alternative to hallucinating, which measurably reduces
fabrication; the example removes ambiguity about what ``reason'' means.

\smallskip
\textbf{And the limit.} Even attempt 3 produces a \emph{drafting aid}, not a risk
model. It has no validated accuracy, no threshold chosen from costs, and no
evaluation against outcomes. For an actual deployed risk score, use
\chref{evaluation}. Knowing which tool the problem calls for is the point.
\end{worked}

\begin{conceptbox}[Chain of Thought]
Adding ``work through this step by step'' measurably improves accuracy on
multi-step reasoning, because the intermediate tokens give the model somewhere to
do the work. It also makes the answer auditable: you can see where the reasoning
went wrong instead of only that it did.
\end{conceptbox}

\section{Retrieval-Augmented Generation}

\begin{definitionbox}[RAG]
Instead of relying on what the model absorbed in training, \emph{retrieve} relevant
passages from your own document store at query time and place them in the prompt.
The model then answers from supplied evidence rather than from memory.
\end{definitionbox}

\dsfig{%
\begin{tikzpicture}[font=\scriptsize,
  b/.style={rounded corners=3pt, draw=#1, fill=#1!9, text=#1!50!black,
            text width=2.35cm, align=center, font=\tiny\bfseries, minimum height=1.0cm},
  arr/.style={-{Stealth[length=2mm]}, gray!60, line width=0.9pt}]

% offline indexing
\node[font=\tiny\bfseries, text=gray!55!black, anchor=west] at (-1.2,1.85) {OFFLINE (once)};
\node[b=primary] (docs) at (0,1.05)  {your documents};
\node[b=secondary] (chunk) at (2.75,1.05) {chunk};
\node[b=secondary] (embed) at (5.5,1.05)  {embed\\{\tiny \chref{nlp}}};
\node[b=purplehl] (vdb) at (8.25,1.05)    {vector\\database};
\foreach \a/\b in {docs/chunk, chunk/embed, embed/vdb}{\draw[arr] (\a) -- (\b);}

% online query
\node[font=\tiny\bfseries, text=gray!55!black, anchor=west] at (-1.2,-0.35) {AT QUERY TIME};
\node[b=primary] (q) at (0,-1.15)     {user question};
\node[b=secondary] (qe) at (2.75,-1.15) {embed the\\question};
\node[b=purplehl] (ret) at (5.5,-1.15)  {retrieve top-$k$\\similar chunks};
\node[b=goldhl] (aug) at (8.25,-1.15)   {question +\\retrieved text};
\node[b=accent] (llm) at (11.0,-1.15)   {LLM};
\node[b=greenhl] (ans) at (13.0,-1.15)  {answer\\+ citations};

\foreach \a/\b in {q/qe, qe/ret, ret/aug, aug/llm, llm/ans}{\draw[arr] (\a) -- (\b);}
\draw[arr, purplehl] (vdb) -- (ret);

\node[anchor=north west, align=left, text width=13.6cm, font=\scriptsize, text=gray!55!black]
     at (-1.2,-2.55)
     {\textbf{What each part fixes:} chunking makes documents small enough to fit the
      context window; embedding makes ``similar in meaning'' computable
      (cosine similarity, \chref{math}); the vector database makes retrieval fast at
      scale; \textbf{citations make the answer checkable} --- which is the single most
      important output of the whole architecture, because it converts an unverifiable
      claim into a verifiable one.};
\end{tikzpicture}}%
{Retrieval-augmented generation. RAG solves three problems at once: the model does not
know your private documents, its training data has a cutoff date, and its claims cannot
otherwise be checked.}{rag}

\section{Evaluating Generative Systems}

\begin{alertbox}[There Is No Accuracy Here]
\chref{evaluation}'s metrics assume one correct answer. A summary has many valid
forms, so precision and recall do not apply. Generative evaluation is therefore
multi-dimensional and partly human --- and a project that has not planned for this
has not planned to know whether it works.
\end{alertbox}

\rowcolors{2}{white}{lightbg}
\begin{longtable}{@{}p{2.9cm}p{4.6cm}p{6.6cm}@{}}
\toprule
\rowcolor{primary!15}
\textbf{Dimension} & \textbf{Question} & \textbf{How to measure} \\
\midrule
\endhead
Groundedness & Is every claim supported by the retrieved evidence? & Automated claim-to-source checking; the key metric for RAG \\
Relevance & Does it answer what was asked? & Human rating, or an LLM judge on a rubric \\
Correctness & Is it factually right? & Gold-answer set; expert review for a sample \\
Safety & Does it refuse what it should refuse? & Red-team prompt suite, run on every release \\
Consistency & Same question, same answer? & Repeat queries at temperature 0 \\
Cost \& latency & Affordable and fast enough? & Tokens per request, p95 response time \\
\bottomrule
\end{longtable}

\begin{pitfall}
\begin{tight}
  \item \textbf{Treating fluency as accuracy.} Confident, well-written and wrong is
        the characteristic failure.
  \item \textbf{Fine-tuning to add facts.} Use RAG; facts change, weights do not.
  \item \textbf{Pasting confidential data into a third-party API.} An acquisition
        and governance question (\chref{ethics}) that must be settled first.
  \item \textbf{No human review for consequential output.} Anything affecting a
        person's grade, employment or security posture needs a human in the loop.
  \item \textbf{Evaluating on the examples you used to write the prompt.} That is
        the leakage of \chref{wrangling} in a new costume.
\end{tight}
\end{pitfall}

%---------------------------------------------------------------------
\begin{fourlenses}{Generative AI}
\lensLA{Draft personalised feedback on assignments, generate practice questions at
a target difficulty, and build a RAG tutor grounded strictly in the module's own
materials so it cannot invent content that is not on the syllabus. \textbf{Firm
limits:} never let a generative model assign a grade unreviewed, and be explicit
with students about where it is used. The fluency of the output makes it far more
persuasive than its reliability warrants.}

\lensPM{Draft status reports from ticket data, summarise long requirement documents,
and extract structured risk registers from meeting notes. RAG over the
organisation's own past project archives gives estimates a reference class it would
otherwise lack. Keep temperature near 0 for anything that must be consistent
between runs --- a status report that changes wording every Monday erodes trust
quickly.}

\lensIOT{Less central here, but real: translating natural-language questions into
telemetry queries, and generating human-readable incident summaries from raw sensor
sequences. RAG over device manuals and past maintenance records puts the right
procedure in a technician's hands at the point of work. Note that LLM inference is
far too heavy for edge devices --- this runs in the cloud, on aggregated data
(\chref{cloudedge}).}

\lensSEC{Two-edged, and the only lens where the technology is used by both sides.
\emph{Defensively:} summarise incidents, explain alerts to junior analysts, RAG over
threat intelligence, draft detection rules. \emph{Offensively:} attackers use the
same models to write flawless phishing in any language, which has removed the
spelling-error signal that classifiers relied on for twenty years. Add prompt
injection to your threat model: if your assistant reads untrusted content, that
content can contain instructions to it.}
\end{fourlenses}

%---------------------------------------------------------------------
\begin{lab}[A Minimal RAG Pipeline]
\begin{lstlisting}[language=Python]
import numpy as np

# --- 1. chunk your documents ------------------------------------------
def chunk(text, size=800, overlap=100):
    """Overlap keeps sentences from being cut across a boundary."""
    return [text[i:i+size] for i in range(0, len(text), size - overlap)]

chunks = [c for doc in documents for c in chunk(doc)]

# --- 2. embed once, offline -------------------------------------------
from sentence_transformers import SentenceTransformer
encoder = SentenceTransformer("all-MiniLM-L6-v2")
index = encoder.encode(chunks, normalize_embeddings=True)   # (n_chunks, 384)

# --- 3. retrieve by cosine similarity (Chapter 4) ---------------------
def retrieve(question, k=4):
    q = encoder.encode([question], normalize_embeddings=True)[0]
    scores = index @ q                       # normalised => dot product IS cosine
    return [(chunks[i], float(scores[i])) for i in np.argsort(scores)[::-1][:k]]

# --- 4. build a grounded prompt ---------------------------------------
def build_prompt(question):
    hits = retrieve(question)
    evidence = "\n\n".join(f"[{i+1}] {c}" for i, (c, _) in enumerate(hits))
    return f"""ROLE: You answer questions about our internal documentation.

CONTEXT (the ONLY permitted source):
{evidence}

QUESTION: {question}

RULES:
- Answer only from the context above.
- Cite the bracketed source number after each claim.
- If the context does not contain the answer, reply exactly:
  "Not covered in the supplied documents."
"""

print(build_prompt("What is our password rotation policy?"))

# --- 5. the evaluation you must not skip ------------------------------
# Build 30-50 question/answer pairs from documents you know. Measure:
#   (a) retrieval hit rate -- was the right chunk in the top k at all?
#   (b) groundedness       -- is every claim traceable to a cited chunk?
#   (c) refusal rate       -- does it correctly refuse unanswerable questions?
# (a) is where most RAG systems actually fail; a perfect model cannot answer
# from evidence it was never given.
\end{lstlisting}
\end{lab}

%---------------------------------------------------------------------
\begin{checkpoint}
\begin{tightnum}
  \item Explain in two sentences why hallucination is a consequence of the
        generation mechanism rather than a bug to be patched.
  \item Your assistant does not know your company's leave policy. Should you
        fine-tune or use RAG? Justify.
  \item You need identical output for the same input every time. What temperature,
        and why?
\end{tightnum}
\end{checkpoint}

%---------------------------------------------------------------------
\begin{chaptersummary}
\begin{tight}
  \item An LLM samples the next token from a learned distribution. Fluency and
        fabrication come from the same mechanism.
  \item Temperature trades determinism against variety and should be set from the
        task.
  \item Adapt in order: prompting, then RAG, then fine-tuning. Fine-tuning teaches
        behaviour, not facts.
  \item A good prompt has role, context, task, format and examples, plus an explicit
        instruction on what to do when evidence is missing.
  \item RAG grounds answers in your own documents and, crucially, makes them
        checkable through citations.
  \item Generative systems have no single correct answer, so evaluate across
        groundedness, relevance, correctness, safety, consistency and cost --- with
        humans in the loop for consequential output.
\end{tight}
\end{chaptersummary}

\begin{reviewq}
  \item \textbf{(MCQ)} An LLM generates text by: (a)~looking up facts in a database
        (b)~predicting the next token repeatedly (c)~searching the web
        (d)~translating from an internal logic.
  \item \textbf{(MCQ)} To make a model answer from your private, frequently-updated
        documents, the appropriate technique is: (a)~pre-training (b)~fine-tuning
        (c)~RAG (d)~raising the temperature.
  \item \textbf{(Short)} Name the five components of a well-formed prompt and say
        what each contributes.
  \item \textbf{(Short)} Explain why citations are the most valuable output of a RAG
        system.
  \item \textbf{(Short)} Give two reasons why generative systems cannot be evaluated
        with precision and recall.
  \item \textbf{(Applied)} Design a RAG assistant that answers student questions
        about a module. Specify the document store, chunking strategy, retrieval
        method, prompt rules, and the three evaluation measures you would report ---
        plus the one failure mode you would monitor most closely.
  \item \textbf{(Applied)} Your security team wants an LLM to summarise alerts. Write
        the threat model: list three ways this could be attacked or could fail
        dangerously, and state a mitigation for each.
\end{reviewq}
